
% !TeX root = ../main.tex
\begin{englishabstract}
    % \footnotetext{*The work was supported by the Foundation (foundation ID).}
    % \astfootnote{The work was supported by the Foundation (foundation ID).}
    % \blindtext

    % \blindtext

    % \blindtext

    % \blindtext
    Three-dimensional reconstruction technology has achieved significant breakthroughs after decades of development and has been widely applied in fields such as industrial inspection, autonomous driving, and cultural heritage preservation, providing high-precision spatial information support for downstream tasks like target recognition and scene understanding. However, traditional image-based three-dimensional reconstruction methods still face limitations in certain specific scenarios, such as logistics sorting and security inspection, where target objects are often covered by packaging, and environmental lighting conditions are restricted or visual occlusion occurs, causing image-based methods to fail. Existing radio frequency signal-based three-dimensional reconstruction methods often rely on complex RF systems with high hardware costs. Additionally, deep learning-based approaches typically depend on large amounts of labeled real data and suffer from high migration costs. Therefore, developing a low-cost, penetrative, and easily transferable three-dimensional reconstruction method holds significant research importance. This thesis addresses the above issues through two main research directions:

    First, to address the failure of existing image-based three-dimensional reconstruction methods in visually occluded scenarios, a novel 3D point cloud reconstruction method based on signal multi-view depth inverse projection and fusion is proposed. This method utilizes only a low-cost single-input single-output commercial ultra-wideband radar to construct a virtual antenna array for signal acquisition of objects. Background reflection reduction and effective range bin localization are employed to obtain denoised target signals. To address the lack of diversity in radar data, virtual antenna offset and phase rotation data augmentation methods are designed. Subsequently, a signal-based multi-view depth reconstruction network extracts three-dimensional structural information from radar signals and generates multi-view depth maps, which are then used for depth map back-projection and point cloud fusion to generate reconstructed point clouds. Experimental results on a self-built dataset show that the Chamfer Distance (CD) error of the reconstructed point cloud is 5.326 cm, with F1 scores of 62.17\% and 76.34\% at 2 cm and 4 cm thresholds, respectively. In fully occluded scenarios, the CD error of the reconstructed point cloud increases by only 0.977 cm compared to non-occluded scenarios, demonstrating the robustness of this method in occluded environments.
    
    Second, to address the reliance of deep learning-based three-dimensional reconstruction methods on large amounts of labeled real data and the high migration costs, a 3D point cloud reconstruction method based on signal domain transfer is proposed. A radar signal simulation algorithm is designed and implemented by detailed modeling of object reflection signals. Additionally, a signal domain adaptation neural network module is proposed to reduce the distribution discrepancy between simulated and real domains, resulting in the construction of a large-scale simulated signal dataset. Transfer learning is then applied, where the signal-based multi-view depth reconstruction network is pre-trained on the simulated signal dataset and fine-tuned on the real signal dataset to enhance data efficiency, training speed, and performance. Under the proposed signal domain transfer-based method, the CD error of the reconstructed point cloud decreases by 0.458 cm, achieving higher reconstruction accuracy. When only 5\% of the training data from the real signal dataset is retained for fine-tuning the signal-based multi-view depth reconstruction network, the CD error of the reconstructed point cloud is 5.175 cm, increasing by only 0.307 cm compared to full-data fine-tuning, significantly improving the data efficiency of the network. Furthermore, the training epochs of the signal-based multi-view depth reconstruction network on the real signal dataset are reduced from 50 to just 8, greatly lowering the training cost.

\englishkeywordstype{IR-UWB Radar; Wireless Sensing; 3D Reconstruction; Transfer Learning}{Application Research}
\end{englishabstract}